% !PRINTED 573
\weberpage{573}

By §\,150, (8), (10), (13) the development of $\tau(u)$ in ascending powers of $u$ begins with a negative power of $u$. In the equation

\weq{4}{\tau(\mu u)\, N = Z}

we arrange the function $Z$ in descending powers of $\tau$, develop both sides in ascending powers of $u$, and compare the coefficients of equally high powers of $u$. We then obtain, for the determination of the coefficients of $Z$, a sequence of linear equations, each following one of which contains only one new one of these coefficients, and from this these coefficients can be calculated successively and rationally. If one now attends to the form of the developments of §\,150, it follows that the coefficients of $Z$ are rational functions of $j(\omega)$ and $\sqrt{\Delta}$.

If one frees $N$ and $Z$ from common factors by the Euclidean algorithm, one comes back to the functions $\Psi, \Phi$, and it follows, if by the class field we understand the totality of the rational functions of $\sqrt{\Delta}$ and $j(\omega)$:

\begin{quote}
3. \gesperrt{The coefficients of the functions $\Phi(\tau)$, $\Psi(\tau)$ in (2) belong to the class field.}
\end{quote}

\begin{center}
§\,154. \textbf{The Division Field.}
\end{center}

If we free the function $\Phi(x)$ (by rational computation) from all factors occurring in it more than once, there remains an integral function $T_\mu(x)$ of $x$ whose coefficients belong to the class field, and the roots of $T_\mu(x)$ are all those different from one another among the quantities

\weq{1}{\tau\left(\frac{\nu}{\mu}\right).}

Two of these quantities

\[
\tau\left(\frac{\nu}{\mu}\right), \quad \tau\left(\frac{\nu'}{\mu}\right)
\]

are equal to one another,

% [sic] In case b) the source also prints a congruence sign before +/-i*nu'/mu; (2) on the next page separates the alternatives by commas. The printed sign is retained.
\[
\begin{array}{rl}
\text{in case (1) a), when } & \dfrac{\nu}{\mu} \equiv \pm\dfrac{\nu'}{\mu}, \\
\text{(1) b), when } & \dfrac{\nu}{\mu} \equiv \pm\dfrac{\nu'}{\mu} \equiv \pm i\dfrac{\nu'}{\mu}, \\
\text{(1) c), when } & \dfrac{\nu}{\mu} \equiv \pm\dfrac{\nu'}{\mu}, \quad \pm\varrho\dfrac{\nu'}{\mu}, \quad \pm\varrho^2\dfrac{\nu'}{\mu},
\end{array}
\]
%CONT>>
